\section{Related Work}
\label{sec:related}

\subsection{Federated Learning Algorithms}

The canonical FedAvg algorithm~\cite{mcmahan2017fedavg} aggregates client
models as a weighted average proportional to local dataset size, but suffers
from \emph{client drift} under non-IID data distributions.
FedProx~\cite{li2020fedprox} adds a proximal regularisation term
$\frac{\mu}{2}\|\mathbf{w}-\mathbf{w}^r\|^2$ that anchors local updates to
the global model, improving convergence under heterogeneity.
FedNova~\cite{wang2020fednova} identifies \emph{objective inconsistency} as
a root cause of client drift and corrects it by normalising each client's
pseudo-gradient by the number of local steps before aggregation.
FedPer~\cite{arivazhagan2019fedper} takes a personalisation perspective:
only the lower network layers (encoder) are aggregated globally, while the
task-specific decoder is kept local, allowing clients to specialise without
sacrificing a shared representation.
Recent aggregation improvements include evidential weight aggregation
(FedEvi~\cite{chen2024fedevi}), which uses Dirichlet-based uncertainty
models to adaptively weight client contributions; and client-contribution
normalisation~\cite{linardos2024fets}, which generalises the FedNova
principle and consistently improves convergence across heterogeneity levels.

\subsection{Federated Learning for Medical Image Analysis}

Sheller \emph{et al.}~\cite{sheller2018federated} demonstrated that FL
models for brain tumour segmentation (BraTS) approach centralised accuracy
while preserving data locality.
NVIDIA FLARE~\cite{nvidia2019clara} provides an open-source infrastructure
for FL in healthcare.
Bercea \emph{et al.}~\cite{bercea2022federated} studied the effect of
non-IID distributions in abdominal CT segmentation and showed that
aggregation strategy selection is critical.
Our work extends these studies by (i) covering a broader set of modalities
and tasks, (ii) comparing four aggregation algorithms with controlled
Dirichlet heterogeneity, and (iii) integrating active learning to address
annotation scarcity.
More recent work has extended FL to foundation-model architectures:
Liu \emph{et al.}~\cite{liu2024fedfms} federate fine-tuning of SAM for
medical segmentation, while Skorupko \emph{et al.}~\cite{skorupko2025fednnunet}
extend nnU-Net with federated fingerprint extraction.
Schutte \emph{et al.}~\cite{schutte2024fedgs} address content heterogeneity
via gradient scaling, and Xiang \emph{et al.}~\cite{xiang2024fedia} handle
annotation incompleteness across clients.
The FeTS 2024 challenge~\cite{linardos2024fets} benchmarks six aggregators
on glioma segmentation, confirming that no single method dominates across
all heterogeneity levels --- motivating the comparative study in this paper.

\subsection{Task-Specific Federated Learning for Ultrasound and Endoscopy}

FL has also been specialised for specific imaging tasks.
For thyroid nodule segmentation in ultrasound, Xiang \emph{et al.}~\cite{xiang2025maunet}
propose a multi-attention UNet federated across multiple ultrasound centres,
achieving SOTA accuracy while preserving data privacy --- directly
complementary to our TN3K experiments.
For polyp segmentation, Pan \emph{et al.}~\cite{pan2025fdgpolyp} apply
Fourier-domain augmentation within FL to improve cross-domain generalisation,
and Tan \emph{et al.}~\cite{tan2025fedsemidg} combine self-supervised
pretraining with adversarial augmentation to reduce label requirements in a
federated setting --- tasks closely related to our FUGC benchmark.
Chen \emph{et al.}~\cite{chen2025pcrfed} demonstrate that contrastive
representation regularisation improves personalised FL for segmentation
across non-IID clients.
In the cardiovascular domain, FedCVD~\cite{zhang2024fedcvd} presents the
first real-world FL benchmark on echocardiogram segmentation, revealing
challenges unique to natural (non-simulated) data heterogeneity.

\subsection{Active Learning for Segmentation}

Active learning~\cite{settles2009active} aims to maximise model accuracy
subject to a labelling budget by selecting the most informative unlabelled
samples.
Entropy-based methods~\cite{settles2009active} query samples with the
highest predictive uncertainty.
BADGE~\cite{ash2020badge} uses gradient embeddings to select a
diverse and uncertain batch.
Coreset~\cite{sener2018coreset} selects samples that minimise the maximum
distance to the existing labelled set in feature space.
The concurrent work most closely related to FedAL is FEAL~\cite{chen2024feal}
(CVPR 2024), which proposes \emph{federated evidential active learning}
using Dirichlet-based epistemic uncertainty to select samples under domain
shift across clients.
Unlike FEAL, our FedAL uses lightweight entropy scoring and is
aggregation-agnostic, making it directly composable with FedAvg, FedProx,
FedNova, and FedPer.
